\subsection{GROUPS OPERATING PROPERLY ON A TOPOLOGICAL SPACE}

DEFINITION 1. Let G be a topological group operating continuously on a topological space X. G is said to operate properly on X if the mapping

\[
\theta : (s, x) \rightarrow (x, s. x) \text{ of } \mathbf{G} \times \mathbf{X} \text{ into } \mathbf{X} \times \mathbf{X}
\]

is proper (Chapter I, § 10, no. 1, Definition 1).

Let \(\Gamma \subset \mathbf{G} \times \mathbf{X} \times \mathbf{X}\) be the graph of the mapping \(\rho: (s, x) \to s.x\) . Since \(\rho\) is continuous, the mapping \(\sigma: (s, x) \to (s, x, s.x)\) is a homeomorphism of \(\mathbf{G} \times \mathbf{X}\) onto \(\Gamma\) , and the composite mapping

\[
\mathbf {G} \times \mathbf {X} \xrightarrow {\sigma} \Gamma \xrightarrow {\mathrm{pr} _ {2 3}} \mathbf {X} \times \mathbf {X}
\]

is just \(\theta\) . Definition 1 is therefore equivalent to saying that the restriction of \(\operatorname{pr}_{23}\) to \(\Gamma\) is a proper mapping of \(\Gamma\) into \(\mathbf{X} \times \mathbf{X}\) .

Theorem 1 of Chapter I, § 10, no. 2 shows that G operates properly on X if and only if the following condition is satisfied:

For each set A filtered by an ultrafilter F, and each mapping

\[
\alpha \rightarrow (s _ {\alpha}, x _ {\alpha})
\]

of A into \(\mathbf{G} \times \mathbf{X}\) , if the mapping \(\alpha \to (s_{\alpha}.x_{\alpha}, x_{\alpha})\) has a limit \((b, a)\) with respect to \(\mathfrak{F}\) , then \(\alpha \to s_{\alpha}\) has a limit \(t \in \mathbf{G}\) with respect to \(\mathfrak{F}\) , such that \(t.a = b\) .

Examples. 1) Let H be a closed subgroup of a topological group G. If G operates properly on X, then so does H, since H × X is closed in G × X (Chapter I, § 10, no. 1, Proposition 5, Corollary 1). If for example we take X = G, with G operating on itself by left translations, then the mapping G × X → X × X is a homeomorphism, hence proper; thus H operates properly on G by left translations.

\begin{enumerate}
\def\labelenumi{\arabic{enumi})}
\setcounter{enumi}{1}
\tightlist
\item
  If G operates properly on X, then it operates properly on every subspace \(X'\) of X which is a union of orbits of points of X (in other words, \(X'\) is saturated with respect to the equivalence relation defined by G). For the inverse image of \(X' \times X'\) in \(G \times X\) is \(G \times X'\) , and we can apply Proposition 3 of Chapter I, § 10, no. 1.
\end{enumerate}

PROPOSITION 1. Let G be a topological group operating continuously on a topological space X, and let K be a quasi-compact subset of G. Then the mapping \(\rho: (s, x) \to s.x\) of \(K \times X\) into \(X\) is proper.

\(\rho\) factorizes into \(\mathbf{K} \times \mathbf{X} \to \mathbf{K} \times \mathbf{X} \xrightarrow{\mathrm{pr}_2} \mathbf{X}\) , where \(\alpha(s, x) = (s, s.x)\) . \(\alpha\) is a homeomorphism, for \(\alpha^{-1}: (s, y) \to (s, s^{-1}.y)\) is continuous. Since \(\mathbf{K}\) is quasi-compact, \(\mathrm{pr}_2\) is proper (Chapter I, § 10, no. 2, Theorem 1, Corollary 5); hence \(\rho\) is proper (Chapter I, § 10, no. 1, Proposition 5).

COROLLARY 1. If A is a closed (resp. compact) subset of X, then K.A is closed in X (resp. compact if X is Hausdorff).

The assertion concerning closed sets follows from Proposition 1 and the fact that a proper mapping is closed (Chapter I, § 10, no. 1, Proposition 1). The assertion concerning compact sets is trivial.

\pandocbounded{\includegraphics[keepaspectratio]{images/part002_62417fce.jpg}}

It should be noted that if L is a compact subset of X, and F a closed subset of G, then F. L is not necessarily closed in X (§ 2, Exercise 29; cf.~no. 5, Corollary to Proposition 12).

COROLLARY 2. If K is a quasi-compact subgroup of a topological group G, then the equivalence relation \(x^{-1}y \in K\) is closed, and the canonical mapping \(\varphi: G \to G/K\) is proper.

The first assertion follows from Corollary 1 applied to G acting by right translations on itself, and the second assertion then follows from Chapter I, § 10, no. 2, Theorem 1.

COROLLARY 3. Let K be a quasi-compact normal subgroup of a topological group G, and let \(\varphi\) be the canonical mapping G → G/K. Then for each closed subgroup A of G the canonical bijection of A/A∩K onto \(\varphi(A)\) is an isomorphism of topological groups.

Since \(x^{-1}y \in K\) is a closed equivalence relation (Corollary 2), the Corollary follows from Chapter I, § 5, no. 2, Proposition 4.

PROPOSITION 2. Let K be a compact group operating continuously on a Hausdorff space X. Then:

\begin{enumerate}
\def\labelenumi{\alph{enumi})}
\item
  K operates properly on X.
\item
  The mapping \((s, x) \to s.x\) of \(\mathbf{K} \times \mathbf{X}\) into \(\mathbf{X}\) is proper.
\item
  The canonical mapping of X onto X/K is proper.
\item
  is a consequence of Proposition 1. As to a), since K is compact, \(\operatorname{pr}_{2}:(s,x)\to x\) is proper (Chapter I, § 10, no. 2, Theorem 1, Corollary 5); hence, X being Hausdorff, \((s,x)\to(x,s.x)\) is proper (Chapter I, § 10, no. 1, Proposition 5, Corollary 3). It remains to prove c). By Corollary 1 to Proposition 1, the canonical mapping \(\varphi:X\to X/K\) is closed. If Z is any topological space and if we make K operate trivially on Z, then K operates continuously on \(X\times Z\) and hence the canonical mapping \(X\times Z\to(X\times Z)/K\) is closed. But \((X\times Z)/K\) can be canonically identified with \((X/K)\times Z\) (§ 2, no. 4, Lemma 2 and Chapter I, § 5, no. 3, Corollary to Proposition 8); hence the canonical mapping \(X\times Z\to(X\times Z)/K\) can be identified with \(\varphi\times1\) , and since it is closed for all Z it follows that \(\varphi\) is proper.
\end{enumerate}

COROLLARY 1. Under the hypotheses of Proposition 2, X is compact (resp. locally compact) if and only if X/K is compact (resp. locally compact).

This follows from the fact that the canonical mapping \(X \rightarrow X/K\) is proper, in view of Proposition 9 of Chapter I, § 10, no. 4.

COROLLARY 2. Let G be a Hausdorff topological group and let K be a compact subgroup of G. Then G is compact (resp. locally compact) if and only if G/K is compact (resp. locally compact).

Apply Corollary 1 to K operating on G by right translations.
% semantic-fragments: legacy-input-seam
\relax{}
